% !TeX root = ../教材题库_全部条目_转写工作稿.tex
% Source: XJ-B1-C06-PR-P227-Y0307
% 第6章 / 6.3 统计图表 / 教材 p.227 / PDF 第234页
% OCR未校对：公式与文字均可能有误；未完成数学TeX转写；禁止直接发给学生
% 原题及相邻原解截图（核对用）：../题目原图/XJ-B1-C06-PR-P227-Y0307-01.webp
% 本文件只在转写完成后作为题干使用；OCR 条目尚不是可靠题干。
\begin{bankitem}{p.227 练习 第2题}{第6章 / 6.3 统计图表 / 教材 p.227 / PDF 第234页}{XJ-B1-C06-PR-P227-Y0307}{OCR未校对：公式与文字均可能有误；未完成数学TeX转写；禁止直接发给学生}
\begin{ocrdraft}
[教材p.227 / PDF第234页]
\par
下列统计图分别摘自国家统计局网站，试说明这些统计图所表达的意义．
\par
2．
\par
充分利用该网站的数据资源就自己感兴趣的话题主动搜集数据，并适当整理，制
\par
作统计图表．
\par
2022年全国居民人均消费支出及其构成
\par
2018一2022年国内游客人次及其增长速度
\par
医疗保健其他用品及服务
\par
2120丿L595儿
\par
亿人次
\par
。=国内游客人次一．一比上年增长％
\par
\&6％
\par
2．4％
\par
食品烟酒
\par
200
\par
教育文化娱乐\textbackslash{}
\par
7481儿
\par
150
\par
60．1
\par
2469丿L
\par
30．5％
\par
6055．4
\par
10．1％
\par
100
\par
交通通信一二
\par
40
\par
2\&8
\par
32．5
\par
3195儿
\par
25．3
\par
13．0％
\par
0
\par
衣着
\par
20
\par
2．
\par
生活用品及服务
\par
1365儿
\par
2．
\par
一50
\par
5．6％
\par
2．
\par
1432兀
\par
0
\par
一100
\par
居住
\par
5．8％
\par
5882兀
\par
2018
\par
2019
\par
2020
\par
2021
\par
2022
\par
24．0％
\par
（第2题）
\end{ocrdraft}
\end{bankitem}
